% !TeX root = ../教材题库_全部条目_转写工作稿.tex
% Source: XJ-B1-C02-RV-P062-Y0994
% 第2章 / 章末复习 复习题二 / 教材 p.62 / PDF 第69页
% OCR未校对：公式与文字均可能有误；未完成数学TeX转写；禁止直接发给学生
% 原题及相邻原解截图（核对用）：../题目原图/XJ-B1-C02-RV-P062-Y0994-01.webp
% 本文件只在转写完成后作为题干使用；OCR 条目尚不是可靠题干。
\begin{bankitem}{复习题二 第12题}{第2章 / 章末复习 复习题二 / 教材 p.62 / PDF 第69页}{XJ-B1-C02-RV-P062-Y0994}{OCR未校对：公式与文字均可能有误；未完成数学TeX转写；禁止直接发给学生}
\begin{ocrdraft}
[教材p.62 / PDF第69页]
\par
12．
\par
设二次函数3，一2一十
\par
4．
\par
若方程3，一0有实根，则实数差的取值范围是
\par
（2）
\par
若不等式3，>0的解集为，则实数差的取值范围是
\par
若不等式3，>0的解集为R，则实数差的取值范围是
\par
（3）
\end{ocrdraft}
\end{bankitem}
